% TEMPLATE-MILC-v3.tex - plantilla maestra UNICA de las guias MILC v3.
%
% La usan las tres familias de guias, cada una con su driver:
%   · Grados 6-11  -> scripts/build-guias-g11.py
%   · Bebras       -> scripts/build-guias-bebras.py
%   · Semillero    -> scripts/build-guias-semillero.py
%
% Antes habia tres copias casi identicas (~400 lineas iguales, 6 distintas):
% cada mejora habia que aplicarla tres veces, y en la practica no se hacia
% (el motor de recursos vivia solo en dos de ellas). Lo que varia entre
% familias esta parametrizado en 6 placeholders que cada driver llena
% (se nombran aqui SIN su sintaxis real: el builder reemplaza por texto plano,
% tambien dentro de los comentarios, y un valor multilinea rompe el preambulo):
%
%   HEADER_IZQ        encabezado izquierdo de pagina
%   HEADER_DER        encabezado derecho de pagina
%   PORTADA_ETIQUETA  rotulo sobre el numero grande (GRADO/MOMENTO/MODULO)
%   PORTADA_SERIE     linea de serie de la portada
%   PORTADA_ID        identificador de la guia en la portada
%   PORTADA_CIRCULO   el circulito de grado (°), o vacio si no aplica
%
% El resto de claves las documenta content/guias/_SCHEMA.md. El script valida
% que no queden placeholders sin reemplazar antes de escribir el archivo final.

\documentclass[11pt,letterpaper]{article}
\usepackage[margin=2.0cm]{geometry}
\usepackage[table]{xcolor}
\usepackage{fontspec}
\usepackage[spanish]{babel}
\usepackage{tikz}
% Librerías TikZ para los diagramas de las guías (bloque `recursos:`):
%  · babel  → babel en español vuelve activos caracteres como «>», y TikZ los usa
%             en sus opciones (>=stealth, ->). Sin esta librería, cualquier
%             diagrama con flechas revienta con «\language@active@arg> has an
%             extra }». Debe ir ANTES de que se dibuje nada.
%  · positioning        → sintaxis «below left=2pt and 3pt».
%  · decorations.pathreplacing → llaves ({brace}) para señalar rangos.
\usetikzlibrary{babel,positioning,decorations.pathreplacing}
\usepackage{graphicx}
% Circuitos: el trabajo de electrónica se hace hoy con micro:bit y sus sensores
% integrados (MakeCode, sin cableado), así que no se necesitan esquemáticos.
% Si algún día entran montajes con protoboard, basta descomentar esta línea:
% los diagramas .tex/.tikz declarados en `recursos:` ya se incrustan solos.
% \usepackage{circuitikz}
\usepackage[most]{tcolorbox}
\usepackage{tabularx}
\usepackage{array}
\usepackage{fancyhdr}
\usepackage{titlesec}
\usepackage[document]{ragged2e}  % bandera a la derecha en todo el cuerpo
\usepackage{enumitem}
\usepackage{microtype}

% Helvetica Neue no trae la flecha U+2192, asi que cada «→» escrito literalmente
% en el YAML se compilaba a NADA, en silencio (462 flechas en 61 guias: rutas del
% tipo «paleta → tipografia → lockup» salian rotas). Mapeamos el caracter a su
% version matematica, que si tiene glifo. Asi el autor puede seguir escribiendo
% «→» comodamente en el YAML.
\usepackage{newunicodechar}
\newunicodechar{→}{\ensuremath{\rightarrow}}
% Mismo problema con los simbolos de marca y vinetas: el «✓» de «una palomita ✓»
% salia como cuadrito vacio en el PDF de todas las guias que lo usaban.
\usepackage{pifont}
\newunicodechar{✓}{\ding{51}}
\newunicodechar{✗}{\ding{55}}
\newunicodechar{✔}{\ding{52}}
\newunicodechar{•}{\textbullet}
\newunicodechar{≥}{\ensuremath{\geq}}
\newunicodechar{≤}{\ensuremath{\leq}}

\setmainfont{Helvetica Neue}
\setsansfont{Helvetica Neue}
\setlength{\parindent}{0pt}
\setlength{\parskip}{6pt}
% Sin particion de palabras: en 6.º-7.º una palabra cortada por guion al final
% de linea («Comuni-cacion») frena la lectura mucho mas que una linea corta.
\hyphenpenalty=10000
\exhyphenpenalty=10000
\pretolerance=2000
\tolerance=3000
\setlength{\emergencystretch}{3em}
\linespread{1.10}
\renewcommand{\arraystretch}{1.25}
\setlist{leftmargin=5mm,itemsep=1mm,topsep=1mm}
\newcolumntype{Y}{>{\RaggedRight\arraybackslash}X}

\definecolor{milcMagenta}{HTML}{E6007E}
\definecolor{milcVino}{HTML}{5A0038}
\definecolor{milcTurquesa}{HTML}{0093A5}
\definecolor{milcVerde}{HTML}{58A923}
\definecolor{milcMostaza}{HTML}{E5B400}
\definecolor{milcOcre}{HTML}{B86600}
\definecolor{milcNegro}{HTML}{241816}
\definecolor{milcGris}{HTML}{F7F7F4}
\definecolor{milcLinea}{HTML}{E8E2DD}
\definecolor{milcGradePrimary}{HTML}{84CC16}
\definecolor{milcGradeDark}{HTML}{4D7C0F}
\definecolor{milcGradeSoft}{HTML}{ECFCCB}
\definecolor{milcGradeAccent}{HTML}{CFE7FF}
\definecolor{milcGradeText}{HTML}{FFFFFF}
\color{milcNegro}

\pagestyle{fancy}
\fancyhf{}
\lhead{\small Tecnología e Informática · Grado 7}
\rhead{\small Guía 16 · MILC v3}
\cfoot{\small\thepage}
\renewcommand{\headrulewidth}{0pt}

\newtcolorbox{titlebox}[1]{
  enhanced,colback=#1,colframe=#1,arc=7mm,boxrule=0pt,
  left=7mm,right=7mm,top=3mm,bottom=3mm,
  before skip=8pt,after skip=8pt,
  fontupper=\sffamily\bfseries\Large\color{white}
}

\newtcolorbox{softbox}[3]{
  enhanced,colback=#2,colframe=#1,borderline west={4pt}{0pt}{#1},
  arc=2mm,boxrule=.4pt,left=4mm,right=4mm,top=3mm,bottom=3mm,
  before skip=6pt,after skip=8pt,title={#3},coltitle=white,
  colbacktitle=#1,fonttitle=\sffamily\bfseries,fontupper=\RaggedRight,
  attach boxed title to top left={xshift=3mm,yshift=-2mm},
  boxed title style={arc=2mm, boxrule=0pt}
}

\newtcolorbox{drawbox}[2]{
  enhanced,colback=white,colframe=#1,arc=4mm,boxrule=1pt,
  left=5mm,right=5mm,top=4mm,bottom=4mm,before skip=8pt,after skip=8pt,
  title={#2},coltitle=white,colbacktitle=#1,fonttitle=\sffamily\bfseries,
  fontupper=\RaggedRight,
  attach boxed title to top left={xshift=4mm,yshift=-2mm},
  boxed title style={arc=3mm, boxrule=0pt}
}

\newtcolorbox{infoband}[1]{
  enhanced,colback=#1!8,colframe=#1!50,arc=2mm,boxrule=.5pt,
  left=4mm,right=4mm,top=2mm,bottom=2mm,before skip=4pt,after skip=4pt,
  fontupper=\small\RaggedRight
}

% Puente narrativo entre secciones (contrato editorial Conéctate).
% Da continuidad: "Acabas de X. Ahora vas a Y porque Z."
% Visualmente: banda izquierda en color del grado, texto en itálica pequeña.
\newtcolorbox{puentebox}{
  enhanced,colback=milcGradeSoft,colframe=milcGradeDark,
  borderline west={3pt}{0pt}{milcGradeDark},
  arc=0mm,boxrule=0pt,left=5mm,right=4mm,top=2.5mm,bottom=2.5mm,
  before skip=4pt,after skip=8pt,
  fontupper=\itshape\small\color{milcNegro}\RaggedRight
}
% La flecha va en modo matematico a proposito: Helvetica Neue NO tiene el glifo
% U+2192, asi que \textrightarrow se compilaba a NADA («Missing character: There
% is no → in font Helvetica Neue») y el puente salia sin flecha. En modo
% matematico la flecha viene de la fuente matematica y si se dibuja.
\newcommand{\puente}[1]{%
  \begin{puentebox}%
    \textcolor{milcGradeDark}{$\rightarrow$}\hspace{2mm}#1%
  \end{puentebox}%
}

\newcommand{\linead}[1]{\noindent\color{milcLinea}\rule{#1}{0.5pt}\color{milcNegro}}
\newcommand{\respuesta}[1]{\par\vspace{2mm}\foreach \n in {1,...,#1}{\noindent\color{milcLinea}\rule{\linewidth}{0.5pt}\par\vspace{5mm}}\color{milcNegro}}
\newcommand{\checkbox}{\textcolor{milcGradeDark}{\fbox{\phantom{x}}}\hspace{5pt}}

\newcommand{\phasecard}[4]{
\begin{tcolorbox}[enhanced,colback=#3,colframe=#2,arc=3mm,boxrule=.5pt,height=3.05cm,valign=center,halign=center,left=1.5mm,right=1.5mm,top=2mm,bottom=2mm]
{\sffamily\bfseries\fontsize{22}{24}\selectfont\textcolor{#2}{#1}}\par
{\sffamily\bfseries\small #4}
\end{tcolorbox}
}

% ── Piezas del rediseno 2026 (legibilidad 6.º-11.º) ───────────────────
% Banda de estacion: reemplaza el titlebox macizo. Titulo a la izquierda,
% dato practico (tiempo/modalidad) a la derecha, en una sola linea.
\newcommand{\milcestacion}[2]{%
\begin{tcolorbox}[enhanced,colback=#1,colframe=#1,arc=2mm,boxrule=0pt,
  left=5mm,right=5mm,top=2.6mm,bottom=2.6mm,before skip=11pt,after skip=7pt]
{\sffamily\bfseries\fontsize{15}{18}\selectfont\color{white}#2}
\end{tcolorbox}}

% Tarjeta de paso numerada: sustituye las tablas «Pilar / Aplicacion», que
% eran formato de consulta docente, no de estudiante. El numero va en un
% circulo sobre el borde para que la secuencia se lea de un vistazo.
\newcommand{\milcpaso}[3]{%
\begin{tcolorbox}[enhanced,colback=white,colframe=milcLinea,arc=1.5mm,boxrule=.9pt,
  left=14mm,right=4mm,top=3mm,bottom=3mm,before skip=4pt,after skip=4pt,
  fontupper=\RaggedRight,
  overlay={\node[anchor=north west,circle,fill=#2,text=white,inner sep=0pt,
    minimum size=7.4mm,font=\sffamily\bfseries\small]
    at ([xshift=3.6mm,yshift=-3.2mm]frame.north west) {#1};}]
#3
\end{tcolorbox}}

% Casilla marcable de tamano util con lapiz (la anterior era \fbox{\phantom{x}},
% de alto variable segun el texto que la seguia).
\newcommand{\milcmarca}{\framebox[3.8mm]{\rule{0pt}{3.4mm}}}

% Fila marcable de lista suelta (checks de la estacion 1).
\newcommand{\milccheckrow}[1]{%
\par\vspace{1.8mm}\noindent
\makebox[6.4mm][l]{\raisebox{-.55mm}{\textcolor{milcNegro}{\milcmarca}}}%
\parbox[t]{\dimexpr\linewidth-6.4mm\relax}{\RaggedRight #1}\par}

% Celda marcable dentro de la rubrica (tabla de autoevaluacion). Misma
% sangria francesa, pero pensada para vivir dentro de una columna X.
\newcommand{\milcopcion}[1]{%
\makebox[5.2mm][l]{\raisebox{-.55mm}{\textcolor{milcNegro}{\milcmarca}}}%
\parbox[t]{\dimexpr\linewidth-5.2mm\relax}{\RaggedRight #1}}

% ── Recursos visuales de la guía ──────────────────────────────────────
% Declarados en el bloque `recursos:` del YAML y emitidos por el builder.
% \guiaFigura   → imagen rasterizada o PDF (foto, carta celeste, montaje).
% \guiaDiagrama → fuente TikZ/circuitikz (.tex) incrustada como vector:
%                 es la vía para circuitos y esquemáticos editables.
% #1 = ruta relativa al .tex · #2 = pie (puede ir vacío)
\newcommand{\guiaFigura}[2]{%
\begin{center}
\includegraphics[width=0.88\linewidth,height=0.38\textheight,keepaspectratio]{#1}\\[1.5mm]
{\footnotesize\itshape\textcolor{milcNegro!75}{#2}}
\end{center}
}

\newcommand{\guiaDiagrama}[2]{%
\begin{center}
\input{#1}\\[1.5mm]
{\footnotesize\itshape\textcolor{milcNegro!75}{#2}}
\end{center}
}

\begin{document}
\thispagestyle{empty}

% ── Portada ───────────────────────────────────────────────────────────
% Antes: un rectangulo de color a pagina completa. Se veia bien en pantalla,
% pero en la fotocopiadora del colegio se comia el toner y salia gris sucio.
% Ahora la banda ocupa el tercio superior y el resto va en blanco.
\begin{tikzpicture}[remember picture,overlay]
  \path[fill=milcGradePrimary]
    (current page.north west) rectangle ([yshift=-10.6cm]current page.north east);

  \node[anchor=north west,text=milcGradeText] at ([xshift=2.0cm,yshift=-1.55cm]current page.north west)
    {\sffamily\bfseries\fontsize{12}{14}\selectfont GRADO};
  \node[anchor=north west,text=milcGradeText,opacity=.85] at ([xshift=2.0cm,yshift=-2.35cm]current page.north west)
    {\sffamily\bfseries\fontsize{8.5}{10}\selectfont SERIE GUÍAS · TECNOLOGÍA E INFORMÁTICA};
  \node[anchor=north east,text=milcGradeText,opacity=.85,align=right] at ([xshift=-2.0cm,yshift=-1.55cm]current page.north east)
    {\sffamily\bfseries\fontsize{9}{12}\selectfont I.E. Sor María Juliana\\Cartago, Valle del Cauca};

  \node[anchor=west,text=milcGradeText] (gradonum) at ([xshift=1.85cm,yshift=-7.1cm]current page.north west)
    {\sffamily\bfseries\fontsize{112}{112}\selectfont 7};
  % El circulito de grado (°) solo aplica a las guias de grado («11°»); en
  % Bebras y Semillero el numero grande es un momento/modulo, no un grado.
  % Se posiciona relativo al nodo `gradonum`, no en coordenadas absolutas.
  \draw[milcGradeText,line width=1.6pt]
    ([xshift=2.0mm,yshift=-4.5mm]gradonum.north east) circle (.15cm);

  \node[anchor=west,text=milcGradeText,text width=11.4cm,align=left] at ([xshift=1.5cm,yshift=0.5cm]gradonum.east)
    {
     {\sffamily\bfseries\fontsize{9}{11}\selectfont\MakeUppercase{Período 2 · Algoritmia y pensamiento computacional}}\\[3mm]
     {\sffamily\bfseries\fontsize{21}{25}\selectfont\setlength{\baselineskip}{27pt}Condicionales ---\\[4pt]el si...\\[4pt]entonces\\[4pt]de los\\[4pt]algoritmos}\\[3.5mm]
     {\sffamily\bfseries\fontsize{15}{18}\selectfont Guía 16-7-TIC}
    };
\end{tikzpicture}

\vspace*{9.4cm}
\vfill

{\sffamily\bfseries\fontsize{8}{10}\selectfont\textcolor{milcGradeDark}{LO QUE TE LLEVAS HOY}}

\begin{tcolorbox}[enhanced,colback=white,colframe=milcNegro,arc=2mm,boxrule=1.6pt,
  left=5mm,right=5mm,top=4mm,bottom=4mm,before skip=3pt,after skip=12pt,
  fontupper=\RaggedRight\fontsize{11}{15.5}\selectfont]
Un \textbf{algoritmo en lenguaje natural} que usa al menos \textbf{3 condicionales} (estructura si-entonces-si no). Tema escogido por ti: \emph{``algoritmo del paraguas''}, \emph{``algoritmo de la calificación''}, \emph{``algoritmo del semáforo''}, etc. Más \textbf{su diagrama de flujo} (aplicando S4: rombos con SÍ/NO etiquetados). Más \textbf{tabla de los 6 operadores de comparación} en el cuaderno.
\end{tcolorbox}

\begin{tcolorbox}[enhanced,colback=milcGris,colframe=milcGris,arc=2mm,boxrule=0pt,
  left=5mm,right=5mm,top=3.5mm,bottom=3.5mm,after skip=12pt]
{\sffamily\bfseries\fontsize{8}{10}\selectfont\textcolor{milcNegro!55}{ESTA GUÍA ES DE}}\\[3mm]
\begin{tabularx}{\linewidth}{X p{3.2cm} p{3.2cm}}
\linead{\linewidth} & \linead{3.0cm} & \linead{3.0cm}\\[-1mm]
{\fontsize{8.5}{10}\selectfont\textcolor{milcNegro!55}{Nombre completo}} &
{\fontsize{8.5}{10}\selectfont\textcolor{milcNegro!55}{Grupo}} &
{\fontsize{8.5}{10}\selectfont\textcolor{milcNegro!55}{Fecha}}\\
\end{tabularx}
\end{tcolorbox}

\vfill

\noindent\textcolor{milcLinea}{\rule{\linewidth}{1pt}}\\[2mm]
{\sffamily\bfseries\fontsize{9.5}{12}\selectfont Autor: Dr. Álvaro Cárdenas Orozco}\\[1mm]
{\sffamily\fontsize{8.5}{11}\selectfont\textcolor{milcNegro!55}{Metodología MILC · apertura ancestral · pensamiento computacional · triángulo Dussel–estoicismo–Floridi}}

\newpage

\section*{Apertura: Condicionales --- el si... entonces de los algoritmos}

\begin{softbox}{milcGradeDark}{milcGradeSoft}{Saber ancestral}
\textbf{Doña Mercedes, la abuela cocinera de Cartago, decidía el almuerzo de cada día con un algoritmo de cabeza.} Cuando llegaba al mercado del barrio, hacía decisiones rápidas: \emph{``¿Hay tomate fresco? Si sí, hago sancocho. Si no, hago arroz con lentejas''}. Después: \emph{``¿Es viernes? Si sí, pescado. Si no, lo que se decidió antes''}. Después: \emph{``¿Tengo presupuesto extra? Si sí, agrego carne. Si no, vegetariano hoy''}. \textbf{Cada decisión seguía una estructura simple:} \emph{``SI tal cosa, ENTONCES tal acción; SI NO, otra acción''}. Así doña Mercedes resolvía cada día el menú \textbf{en 2 minutos en el mercado}, sin paralizarse. Esa estructura --- \emph{si... entonces... si no} --- es exactamente el \textbf{condicional} en programación. Las abuelas cocineras, los albañiles que decidían si un mes era invierno o no, los pescadores que escogían si salir al río o no según el clima: todos usaban condicionales sin saberlo.
\end{softbox}

\begin{softbox}{milcTurquesa}{milcGris}{Saber contemporáneo}
Un \textbf{condicional} (o \emph{estructura selectiva}) es un componente fundamental de la programación que permite \textbf{ejecutar una acción solo SI se cumple cierta condición}. La estructura básica es: \emph{SI [condición] ENTONCES [acción 1] SI NO [acción 2]}. En inglés: \emph{IF [condition] THEN [action 1] ELSE [action 2]}. La \emph{condición} es una pregunta que tiene respuesta verdadero o falso (booleano). Si es verdadera, se ejecuta la acción 1; si es falsa, la acción 2. \textbf{Operadores de comparación} (las preguntas básicas que se hacen en los condicionales): \emph{= (igual a)}, \emph{≠ (distinto a)}, \emph{> (mayor que)}, \emph{< (menor que)}, \emph{≥ (mayor o igual)}, \emph{≤ (menor o igual)}. Combinados con variables (las cajas con etiqueta de la S5), los condicionales permiten que un algoritmo \emph{tome decisiones según los valores}. Esto es lo que hace que los programas se sientan \emph{inteligentes}: no siguen siempre el mismo camino, escogen según los datos.
\end{softbox}

\begin{softbox}{milcMostaza}{milcGris}{Pregunta-puente}
\textit{Cada mañana tu mamá te dice: \emph{``si llueve, llevas paraguas. Si no, no.''} ¿\textbf{Cuántas decisiones de este tipo} tomas tú al día sin darte cuenta? ¿5? ¿20? ¿100?}
\end{softbox}

\begin{softbox}{milcVerde}{milcGris}{Saber y saber hacer}
Al terminar la clase: (1) podrás \textbf{identificar} la estructura si-entonces-si no; (2) sabrás \textbf{aplicar} los 6 operadores de comparación; (3) podrás \textbf{evaluar} si un condicional está bien escrito; (4) habrás \textbf{creado} un algoritmo con 3 condicionales + diagrama.
\end{softbox}



\small
\begin{tabularx}{\textwidth}{>{\bfseries}p{3.0cm}Y}
\rowcolor{milcGradePrimary!12}Autor & Dr. Álvaro Cárdenas Orozco\\
DBA & Construye algoritmos y diagramas de flujo para resolver problemas paso a paso (MEN, T\&I 7°)\\
Referentes & Saber ancestral del tejedor · Pensamiento computacional · Scratch\\
Producto final & Un \textbf{algoritmo en lenguaje natural} que usa al menos \textbf{3 condicionales} (estructura si-entonces-si no). Tema escogido por ti: \emph{``algoritmo del paraguas''}, \emph{``algoritmo de la calificación''}, \emph{``algoritmo del semáforo''}, etc. Más \textbf{su diagrama de flujo} (aplicando S4: rombos con SÍ/NO etiquetados). Más \textbf{tabla de los 6 operadores de comparación} en el cuaderno.\\
\end{tabularx}
\normalsize

\puente{Hoy haces 4 cosas: identificas decisiones en situaciones reales, aprendes la estructura + operadores, escribes algoritmo con 3 condicionales, dibujas el diagrama.}

\milcestacion{milcGradeDark}{Ruta MILC: así vamos a aprender}
\begin{tabularx}{\textwidth}{>{\centering\arraybackslash}X>{\centering\arraybackslash}X>{\centering\arraybackslash}X>{\centering\arraybackslash}X}
\phasecard{1}{milcVerde}{milcGris}{Escucha\\[-1mm]\small Observo, escucho y pregunto.} &
\phasecard{2}{milcTurquesa}{milcGris}{Entender\\[-1mm]\small Sistematizo y ordeno.} &
\phasecard{3}{milcMagenta}{milcGris}{Hacer\\[-1mm]\small Construyo y aplico.} &
\phasecard{4}{milcVino}{milcGris}{Cierre\\[-1mm]\small Evalúo y reflexiono.}
\end{tabularx}


\puente{Empezamos identificando decisiones que ya tomas hoy.}

\milcestacion{milcVerde}{Estación 1 de 3 · Escucha}

\begin{softbox}{milcVerde}{milcGris}{Escena para escuchar}
\textbf{Actividad 1 · IDENTIFICA --- Decisiones de tu día} (10 min · individual). En el cuaderno, anota \textbf{6 decisiones} que tomas en un día normal (desde que te despiertas hasta que duermes). Para cada una, escríbela en formato condicional: \emph{``SI [condición], ENTONCES [acción 1]; SI NO, [acción 2]''}. Ejemplos: \emph{``SI llueve, ENTONCES llevo paraguas; SI NO, no lo llevo''}. \emph{``SI hay desayuno en la mesa, ENTONCES desayuno; SI NO, voy directo al colegio''}. \emph{``SI termino la tarea, ENTONCES puedo jugar; SI NO, sigo trabajando''}. Lista 6 reales tuyas.
\end{softbox}

{\sffamily\bfseries\fontsize{8}{10}\selectfont\textcolor{milcNegro!55}{MARCA CUANDO LO HAYAS HECHO}}
\milccheckrow{Pensé en mi día normal.}
\milccheckrow{Anoté 6 decisiones reales.}
\milccheckrow{Las formulé en estructura SI-ENTONCES-SI NO.}

\begin{infoband}{milcVerde}


\textbf{Lo que va al cuaderno (Actividad 1 · IDENTIFICA).} Título: «Actividad 1 --- 6 decisiones de mi día». \textbf{Formato:} lista numerada del 1 al 6 con estructura SI-ENTONCES-SI NO. \textbf{Extensión:} 6 líneas. \textbf{Sabes que terminaste cuando} reconoces que tu día está lleno de condicionales.
\end{infoband}

\puente{Después aprendes la estructura formal + los 6 operadores.}

\milcestacion{milcTurquesa}{Estación 2 de 3 · Entender}

\begin{softbox}{milcTurquesa}{milcGris}{Lo que vas a entender}
Lo que hiciste son \textbf{condicionales de la vida real}. Tomas más de 100 decisiones así al día sin darte cuenta. Ahora aprende la estructura formal + los 6 operadores que se usan para comparar.
\end{softbox}

\milcpaso{1}{milcTurquesa}{\textbf{La estructura del condicional simple.} En lenguaje natural: \emph{SI [condición] ENTONCES [acción]}. Si la condición no se cumple, simplemente no se ejecuta la acción. Ejemplo: \emph{``SI llueve, ENTONCES llevo paraguas''}. \textbf{Estructura del condicional doble (si-no):} \emph{SI [condición] ENTONCES [acción 1] SI NO [acción 2]}. Aquí siempre se ejecuta alguna de las 2 acciones, según la condición. Ejemplo: \emph{``SI llueve, ENTONCES llevo paraguas; SI NO, no lo llevo''}.}
\milcpaso{2}{milcTurquesa}{\textbf{Los 6 operadores de comparación.} Para escribir condiciones, usamos estos símbolos: \textbf{(1) = igual a:} \texttt{edad = 13} ¿es la edad exactamente 13? \textbf{(2) $\neq$ distinto a:} \texttt{nombre $\neq$ ``Pedro''} ¿el nombre no es Pedro? \textbf{(3) > mayor que:} \texttt{precio > 100} ¿el precio es mayor que 100? \textbf{(4) < menor que:} \texttt{altura < 1.50} ¿la altura es menor que 1.50? \textbf{(5) $\geq$ mayor o igual:} \texttt{nota $\geq$ 3.5} ¿la nota es mayor o igual a 3.5 (es decir, aprueba)? \textbf{(6) $\leq$ menor o igual:} \texttt{edad $\leq$ 12} ¿la edad es menor o igual a 12 (es decir, niño)?.}
\milcpaso{3}{milcTurquesa}{\textbf{Condicionales anidados (uno dentro de otro).} Cuando una decisión depende de la anterior, se anidan. Ejemplo del semáforo: \emph{SI luz = ``verde'' ENTONCES seguir. SI NO SI luz = ``amarillo'' ENTONCES disminuir velocidad. SI NO ENTONCES detenerse}. Esto se escribe escalonado para mayor claridad. \textbf{Pista:} no anides más de 3 condicionales seguidos --- se vuelve enredado. Si necesitas más, probablemente es mejor descomponer en sub-algoritmos.}
\milcpaso{4}{milcTurquesa}{\textbf{Conexión con diagramas de flujo (S4).} Cada condicional en el algoritmo se representa con un \textbf{rombo de decisión}: adentro va la condición (pregunta), dos flechas salen con SÍ y NO. \textbf{Cómo dibujarlo:} (a) rombo con la pregunta dentro; (b) flecha SÍ que lleva al rectángulo de la acción 1; (c) flecha NO que lleva al rectángulo de la acción 2; (d) ambas flechas convergen al siguiente paso común. Si tienes 3 condicionales en tu algoritmo, tu diagrama tendrá 3 rombos.}

\begin{drawbox}{milcTurquesa}{Condicionales + 6 operadores}
\textbf{Estructura simple:} SI [condición] ENTONCES [acción]. \\[1mm]
\textbf{Estructura doble:} SI [condición] ENTONCES [acción 1] SI NO [acción 2]. \\[1mm]
\textbf{6 operadores:} \\[1mm]
= igual · $\neq$ distinto · > mayor · < menor · $\geq$ mayor o igual · $\leq$ menor o igual. \\[1mm]
\textbf{En diagrama:} cada condicional es un rombo con SÍ/NO.
\end{drawbox}

\begin{softbox}{milcMostaza}{milcGris}{Errores comunes y su corrección}
\textbf{Errores frecuentes con condicionales.} (1) \emph{Olvidar el SI NO:} solo escribir \emph{``SI llueve, llevo paraguas''} y no decir qué hacer si no llueve. A veces está bien (no se ejecuta nada), pero a veces dejas un hueco. (2) \emph{Confundir = con ==:} en algunos lenguajes el = es asignación y == es comparación. En lenguaje natural y Scratch, = sirve para los 2. (3) \emph{Condicionales sin sentido lógico:} \emph{``SI 5 > 10''} es siempre falso. Asegúrate de que la condición pueda ser verdadera o falsa según la situación. (4) \emph{Anidar demasiado:} 5 condicionales seguidos se vuelven ilegibles. Mejor descomponer. (5) \emph{Olvidar que las cadenas de texto necesitan comillas en las comparaciones:} \texttt{nombre = ``María''}, no \texttt{nombre = María}.
\end{softbox}

\begin{infoband}{milcTurquesa}


\textbf{Lo que va al cuaderno (Actividad 2 · EXPLICA).} Título: «Actividad 2 --- Estructura + 6 operadores». \textbf{Formato:} bloque A con las 2 estructuras (simple y doble) con 1 ejemplo cada una. Bloque B con tabla de 6 filas: operador + significado + ejemplo. \textbf{Extensión:} 10 líneas. \textbf{Sabes que terminaste cuando} puedes leer un condicional en código y traducirlo a lenguaje natural.
\end{infoband}

\puente{Con todo claro, escribes algoritmo + diagrama de flujo.}

\milcestacion{milcMagenta}{Estación 3 de 3 · Hacer}

\begin{softbox}{milcMagenta}{milcGris}{Cómo vas a construir tu producto}
\textbf{Actividad 3 · APLICA --- Algoritmo con 3 condicionales + diagrama} (30 min · individual). Vas a escribir un algoritmo de tu vida que use al menos 3 condicionales, y dibujarlo en diagrama de flujo.
\end{softbox}

\milcpaso{1}{milcMagenta}{\textbf{Paso 1 --- Escoge el tema del algoritmo.} En el cuaderno escribe: \emph{``Mi algoritmo: \_\_\_\_\_''}. Tiene que ser algo que requiera al menos 3 decisiones. Sugerencias: \emph{``Algoritmo del paraguas''} (¿llueve? ¿hace frío? ¿voy a salir mucho?), \emph{``Algoritmo de la calificación''} (¿saca más de 4.5? excelente; ¿más de 3.5? aprueba; ¿menos? falla), \emph{``Algoritmo del fin de semana''} (¿es viernes? ¿es sábado? ¿hay tareas pendientes?), \emph{``Algoritmo de qué ropa ponerme''} (¿es para colegio? ¿hace frío? ¿llueve?).}
\milcpaso{2}{milcMagenta}{\textbf{Paso 2 --- Lista las variables que necesitas.} Si tu algoritmo usa información que cambia, son variables (S5). Ejemplo del paraguas: \texttt{estaLloviendo} (booleano), \texttt{vaASalir} (booleano), \texttt{horasFuera} (número). Lista las variables al inicio del algoritmo.}
\milcpaso{3}{milcMagenta}{\textbf{Paso 3 --- Escribe el algoritmo en lenguaje natural.} Estructura: \emph{INICIO. (entrada de variables). PASO 1 con un condicional. PASO 2 con condicional. PASO 3 con condicional. SALIDA. FIN}. Ejemplo del paraguas: \emph{``INICIO. Variables: estaLloviendo, vaASalir, horasFuera. PASO 1: SI estaLloviendo = verdadero ENTONCES llevarParaguas = verdadero. SI NO, llevarParaguas = falso. PASO 2: SI vaASalir = falso ENTONCES llevarParaguas = falso (no necesito si no salgo). PASO 3: SI horasFuera > 5 ENTONCES llevarParaguasGrande = verdadero. SI NO, llevarParaguasGrande = falso. SALIDA: Decisión final. FIN.''}.}
\milcpaso{4}{milcMagenta}{\textbf{Paso 4 --- Dibuja el diagrama de flujo.} Aplica lo aprendido en S4: óvalo INICIO → procesos → rombos para cada condicional → óvalo FIN. Tu diagrama tendrá \textbf{3 rombos} (uno por condicional). Cada rombo con sus 2 flechas etiquetadas SÍ/NO.}
\milcpaso{5}{milcMagenta}{\textbf{Paso 5 --- Tabla de operadores.} En una esquina del cuaderno, dibuja una tabla con los 6 operadores: =, ≠, >, <, ≥, ≤. Al lado de cada uno, el significado en español. Marca cuáles usaste en tu algoritmo.}
\milcpaso{6}{milcMagenta}{\textbf{Paso 6 --- Prueba mental + reflexión.} Imagina 2 valores distintos para tus variables (\emph{caso 1: lloviendo y voy a salir; caso 2: soleado y no salgo}). Ejecuta el algoritmo mentalmente para los 2 casos y verifica que da resultados correctos. Al final escribe: \emph{``El condicional que más me costó escribir fue \_\_\_. El operador que más usé fue \_\_\_''}. Firma con nombre y fecha.}

\begin{drawbox}{milcMagenta}{Antes de cerrar la S6}
\checkbox Mi algoritmo tiene tema claro y útil. \\[1mm]
\checkbox Listé las variables al inicio. \\[1mm]
\checkbox Hay mínimo 3 condicionales con estructura SI-ENTONCES-SI NO. \\[1mm]
\checkbox El diagrama tiene 3 rombos con SÍ/NO etiquetados. \\[1mm]
\checkbox La tabla de los 6 operadores está al lado. \\[1mm]
\checkbox Probé el algoritmo mentalmente con valores distintos.
\end{drawbox}

\begin{softbox}{milcMagenta}{milcGris}{Plantilla de trabajo}
\textbf{Estructura.} \textit{Paso 1:} tema del algoritmo. \textit{Paso 2:} variables. \textit{Paso 3:} algoritmo con 3 condicionales (SI-ENTONCES-SI NO). \textit{Paso 4:} diagrama de flujo. \textit{Paso 5:} tabla de 6 operadores. \textit{Paso 6:} prueba mental + reflexión.
\end{softbox}

\begin{infoband}{milcMagenta}


\textbf{Lo que va al cuaderno (Actividad 3 · APLICA).} Título: «Actividad 3 --- Algoritmo con 3 condicionales». \textbf{Formato:} variables + algoritmo + diagrama + tabla de operadores + prueba mental. \textbf{Extensión:} 1 página y media. \textbf{Sabes que terminaste cuando} tu algoritmo da resultados correctos para distintos valores de entrada.
\end{infoband}

\puente{El producto es algoritmo + diagrama + tabla de operadores.}

\milcestacion{milcMostaza}{Producto de la sesión}
\textbf{Algoritmo con 3 condicionales + diagrama de flujo · 6 operadores aplicados}.
\begin{softbox}{milcMostaza}{milcGris}{Criterios de calidad}
(1) El algoritmo tiene mínimo 3 condicionales con estructura completa. (2) Las variables están listadas al inicio. (3) El diagrama de flujo tiene 3 rombos con SÍ/NO etiquetados. (4) La tabla de operadores está completa. (5) Se hizo prueba mental con valores distintos.
\end{softbox}

\puente{Antes de salir, verifica que cada condicional tiene su SI NO.}

\milcestacion{milcVino}{Evaluación liberadora · cinco dimensiones}

\begin{softbox}{milcVino}{milcGris}{Sentido de la evaluación}
Evaluar no es solo revisar una nota. Es reconocer qué aprendí, qué cambió en mí, qué emociones manejé, cómo me relaciono con la ciudad y la comunidad, y qué le dejaré a quien viene detrás.
\end{softbox}

\textbf{Mi reflexión liberadora en cinco dimensiones.} Completa cada frase iniciada:

\small
\begin{tabularx}{\textwidth}{>{\bfseries\RaggedRight\arraybackslash}p{3.4cm}Y>{\RaggedRight\arraybackslash}p{4.6cm}}
\rowcolor{milcVino}\color{white}Dimensión & \color{white}\bfseries Mi respuesta (frame starter) & \color{white}\bfseries Algo que descubrí\\
1. Personal & Lo que más cambió en mí fue \linead{6.5cm} & \linead{4.2cm}\\
2. Emocional & La emoción más fuerte fue \linead{2.5cm} y la regulé \linead{3cm} & \linead{4.2cm}\\
3. Ciudadana & En democracia esto importa porque \linead{6cm} & \linead{4.2cm}\\
4. Local (Cartago) & En mi barrio sería útil para \linead{6.5cm} & \linead{4.2cm}\\
5. Intergeneracional & A mi abuela/abuelo le diría \linead{6.5cm} & \linead{4.2cm}\\
\end{tabularx}
\normalsize

\begin{infoband}{milcVino}
\textbf{Para la próxima sustentación.} Vuelve a esta tabla en seis meses. Verás que las respuestas cambian — no porque lo aprendido cambie, sino porque tú cambias. La evaluación liberadora es una conversación contigo a través del tiempo.
\end{infoband}


\puente{Cerramos con 3 voces sobre por qué tomar decisiones con criterio es virtud universal.}

\milcestacion{milcGradeDark}{Triángulo de pensamiento · cierre filosófico}

\begin{softbox}{milcVino}{milcGris}{Dussel · lente del nosotros}
\textit{``Toda decisión es un condicional. SI quiero ser justo, ENTONCES debo escuchar. La vida ética es algoritmo aplicado.''}

Las decisiones éticas también son condicionales: \emph{``SI mi amigo necesita ayuda, ENTONCES la ofrezco''}, \emph{``SI veo una injusticia, ENTONCES no callo''}. Aprender a pensar en condicionales no es solo programación: \textbf{es aprender a estructurar tu vida moral}. Las abuelas cocineras tomaban decisiones técnicas (ingredientes) y éticas (¿alcanza para los vecinos?) con la misma estructura mental. \emph{El condicional es la forma básica del pensamiento responsable}.

\textbf{¿Cuáles son las condicionales éticas que rigen mis decisiones diarias? ¿Las he hecho explícitas?}
\end{softbox}
\linead{\linewidth}\\[1mm]
\linead{\linewidth}

\begin{softbox}{milcMostaza}{milcGris}{Estoicismo · Marco Aurelio · cuidado interior}
\textit{``El sabio no decide por impulso. Pregunta: SI hago esto, ¿qué pasa? SI no, ¿qué pasa? Y elige con criterio.''}

Marco Aurelio gobernaba un imperio. \textbf{Cada decisión la pensaba como condicional}: \emph{``SI envío estas tropas a Galia, ENTONCES protejo el norte pero debilito el sur. SI NO, conservo el balance''}. La habilidad de \emph{pensar en condicionales} antes de actuar es lo que separa al estoico del impulsivo. Tú con tus algoritmos entrenas esa virtud antigua en lenguaje moderno.

\textbf{Antes de tomar decisiones importantes, ¿hago el ejercicio mental del condicional (SI esto, entonces tal; SI no, tal otro)?}
\end{softbox}
\linead{\linewidth}\\[1mm]
\linead{\linewidth}

\begin{softbox}{milcTurquesa}{milcGris}{Floridi · lente de la infoesfera}
\textit{``Los algoritmos modernos toman decisiones por nosotros: qué video ver, qué amigo recomendar. Entender los condicionales es entender cómo deciden esos algoritmos.''}

Cuando TikTok te recomienda un video, está ejecutando algoritmos con \textbf{millones de condicionales} (SI le gustaron videos similares, ENTONCES recomendar este). Cuando un banco aprueba un crédito, también: condicionales sobre tu historial. \emph{Entender qué son los condicionales te abre la caja negra de los algoritmos que rigen tu vida}. No para programar a Big Tech, sino para no ser sujeto pasivo.

\textbf{¿Qué algoritmos invisibles me afectan cada día? ¿Sé al menos que existen?}
\end{softbox}
\linead{\linewidth}\\[1mm]
\linead{\linewidth}

\begin{infoband}{milcNegro}
\textbf{Nota para el docente.} Las tres formulaciones del triángulo son la síntesis del autor de la guía, no citas textuales de los pensadores. Si vas a citarlos en clase, remite a la obra original.
\end{infoband}

\milcestacion{milcVino}{Mi autoevaluación · marca una casilla por fila}

{\small
\setlength{\extrarowheight}{2.2pt}
\begin{tabularx}{\textwidth}{>{\RaggedRight\arraybackslash\bfseries}p{3.0cm}YYY}
\rowcolor{milcVino}\color{white}\bfseries Criterio & \color{white}\bfseries Lo logré & \color{white}\bfseries En proceso & \color{white}\bfseries Necesito apoyo\\[.6mm]
Comprensión del tema & \milcopcion{Explico las ideas con mis palabras.} & \milcopcion{Reconozco las ideas, falta precisión.} & \milcopcion{Necesito repasar lo básico.}\\[.8mm]
\rowcolor{milcGris}Pensamiento computacional & \milcopcion{Usé los pilares al armar mi producto.} & \milcopcion{Usé algunos pilares.} & \milcopcion{Necesito releer la estación 2.}\\[.8mm]
Aplicación y producto & \milcopcion{Mi producto cumple los criterios.} & \milcopcion{Está armado, con ajustes pendientes.} & \milcopcion{Aún está en construcción.}\\[.8mm]
\rowcolor{milcGris}Reflexión 5 dimensiones & \milcopcion{Respondí cada una con honestidad.} & \milcopcion{Respondí algunas con detalle.} & \milcopcion{Mis respuestas son muy generales.}\\[.8mm]
Triángulo de pensamiento & \milcopcion{Conecté las 3 voces con mi trabajo.} & \milcopcion{Conecté al menos una.} & \milcopcion{No las relaciono con mi proceso.}\\[.8mm]
\end{tabularx}}

\vspace{3mm}
\textbf{Hoy entendí que:}\\[1mm]
\linead{\linewidth}\\[1mm]
\linead{\linewidth}

\textbf{Mi compromiso para la próxima sesión es:}\\[1mm]
\linead{\linewidth}\\[1mm]
\linead{\linewidth}

\begin{softbox}{milcMostaza}{milcGris}{Cita de despedida · Marco Aurelio}
\textit{``El obstáculo en el camino se vuelve el camino. Lo que estorba al camino se vuelve camino.''}\\[1mm]
Cada error de hoy, bien observado, se convierte en pieza del camino que pisarás mañana.
\end{softbox}

\small
\begin{tabularx}{\textwidth}{>{\bfseries}p{3.6cm}Y}
\rowcolor{milcGradeDark!12}Nombre completo: & \linead{10cm}\\
Fecha: & \linead{4cm} \quad Curso/grupo: \linead{4cm}\\
Mi firma: & \linead{9cm}\\
\end{tabularx}
\normalsize

\begin{infoband}{milcGradeDark}
\textbf{Para la próxima sesión.} Esta semana, anota 3 condicionales reales de tu vida (decisiones que tomaste con la estructura SI-ENTONCES). La próxima clase aprendes \textbf{bucles}: la otra herramienta fundamental de programación, que permite repetir acciones inteligentemente.
\end{infoband}

\end{document}
